\documentclass[../../../include/open-logic-section]{subfiles}

\begin{document}

\olfileid[hi]{sth}{spine}{valpha}

\olsection{चरणों को $V_\alpha$ के रूप में परिभाषित करना}

\olref[sth][ordinals][]{chap} में हमने सुव्यवस्थाएँ और (वॉन नोयमान)
क्रमसूचक परिभाषित किए। इस अध्याय में हम इनका उपयोग स्वयं समुच्चय पदानुक्रम
को अभिलक्षित करने में करेंगे। इसके लिए याद करें कि
\olref[sth][ordinals][opps]{sec} में हमने उत्तरवर्ती और सीमा क्रमसूचकों की
धारणा परिभाषित की थी। निम्न परिभाषा में हम इन्हीं विचारों का उपयोग करते हैं:

\begin{defn}\ollabel{defValphas}
	\begin{align*}
	V_\emptyset &\defis \emptyset\\
	V_{\ordsucc{\alpha}} &\defis \Pow{V_\alpha} & & 
	\text{किसी भी क्रमसूचक }\alpha\text{ के लिए}\\
	V_{\alpha} &\defis \bigcup_{\gamma < \alpha} V_\gamma & & 
	\text{जब }\alpha\text{ सीमा क्रमसूचक हो}
\end{align*}
\end{defn}
\noindent
यह क्रमसूचकों पर \emph{परिमितेतर पुनरावर्तन} द्वारा परिभाषा होगी। इस संबंध
में हमें इसकी तुलना प्राकृतिक संख्याओं पर फलनों की पुनरावर्ती परिभाषाओं से
करनी चाहिए।\footnote{\olref[sfr][infinite][dedekind]{sec} में योग, गुणन और
घातांक की परिभाषाओं से तुलना करें।} प्राकृतिक संख्याओं की तरह हम आधार मामला
और उत्तरवर्ती मामले परिभाषित करते हैं; किंतु क्रमसूचकों के साथ हमें
\emph{सीमा} मामलों का व्यवहार भी बताना होगा।

$V_\alpha$ की यह परिभाषा एक महत्त्वपूर्ण पड़ाव होगी। हमने समुच्चयों के अपने
पदानुक्रम को अनौपचारिक रूप से \emph{चरणों} में समुच्चय बनाने के रूप में
प्रेरित किया है। वास्तव में $V_\alpha$ ठीक वही चरण हैं। किंतु महत्त्वपूर्ण
बात यह है कि यह चरणों का \emph{आंतरिक} अभिलक्षण है। सिद्धांत का कोई संभव
\emph{प्रतिरूप} सुझाने के बजाय हमने अपने समुच्चय सिद्धांत के \emph{भीतर}
चरण परिभाषित किए होंगे।

\end{document}
